\section{Module \texttt{window}}

Ce module crée la fenêtre principale GTK et applique les propriétés de base (titre, taille, redimensionnement, modalité et style).

\subsection{Macros de Configuration}

\begin{docbox}{Macro \texttt{WINDOW\_CONFIG\_DEFAULT}}
\begin{lstlisting}[language=C]
#define WINDOW_CONFIG_DEFAULT ((window_config){ \
    .title = NULL, \
    .icon_name = NULL, \
    .widget_name = NULL, \
    .background_color = NULL, \
    .background_image_path = NULL, \
    .default_width = -1, \
    .default_height = -1, \
    .min_width = 0, \
    .min_height = 0, \
    .max_width = 0, \
    .max_height = 0, \
    .resizable = true, \
    .decorated = true, \
    .modal = false, \
    .maximized = false, \
    .present_on_create = false, \
    .style = WIDGET_STYLE_CONFIG_DEFAULT, \
})
\end{lstlisting}
\end{docbox}

Cette macro fournit une configuration sûre avec des pointeurs initialisés à \texttt{NULL} et des booléens cohérents. Elle évite de lire des champs indéfinis au moment de la création de la fenêtre. En pratique, on part de cette base puis on surcharge uniquement les options nécessaires.

Cette initialisation par défaut est importante en C: sans valeurs explicites, certains champs peuvent contenir des données aléatoires en mémoire. En utilisant cette macro, on réduit les comportements non déterministes et les erreurs difficiles à reproduire.

\subsection{Fonctions API}

\begin{lstlisting}[language=C]
GtkWidget *create_window(GtkApplication *app, const window_config *config) {
    GtkWidget *window = gtk_application_window_new(app);
    if (config == NULL) {
        return window;
    }

    if (config->title) {
        gtk_window_set_title(GTK_WINDOW(window), config->title);
    }
    if (config->icon_name) {
        gtk_window_set_icon_name(GTK_WINDOW(window), config->icon_name);
    }

    gtk_window_set_default_size(GTK_WINDOW(window), config->default_width, config->default_height);
    gtk_window_set_resizable(GTK_WINDOW(window), config->resizable);
    gtk_window_set_decorated(GTK_WINDOW(window), config->decorated);
    gtk_window_set_modal(GTK_WINDOW(window), config->modal);

    apply_window_background(config, window);
    apply_widget_style(window, &config->style);

    if (config->maximized) {
        gtk_window_maximize(GTK_WINDOW(window));
    }
    if (config->present_on_create) {
        gtk_window_present(GTK_WINDOW(window));
    }

    return window;
}
\end{lstlisting}

La fonction applique une stratégie défensive: toujours retourner un widget valide, puis enrichir selon la configuration. Cela simplifie l'appelant et réduit les cas d'erreur côté application.

Le rôle de \texttt{create\_window} est aussi d'unifier la politique d'initialisation de la fenêtre: dimensions, état de présentation et style CSS sont traités au même endroit. Cette centralisation améliore la maintenabilité et évite de disperser des appels GTK critiques dans plusieurs fichiers.

\newpage
